% !TeX root = ../QFT_main.tex
%=========================================================
%=========================================================
\chapter{Non-perturbative aspects of QFT}
%========================================================
%=========================================================



%--------------------------------------------------------
%=========================================================
\section{Kallén-Lehmann spectral representation}
%=========================================================
%------------------------------------------------------------



\begin{mytheorem}[The states of a theory (Matheus 2.1)]
In order to discuss the Kallén-Lehmann representation, we need to understand the states of a quantum field theory.
We can have essentially three types of states, the first two being essentially the same,
\begin{itemize}
    \item \textbf{Single-particle states.} These are states that correspond to a single excitation of the field, typically associated with a pole in the propagator of the theory.
    These have a well-defined mass and momentum, and are reasonably stable to be identified as particles.
    \item \textbf{Resonances.} These are states that correspond to a single excitation of the field, but are not stable enough to be identified as particles, and thus correspond to a resonance in the propagator of the theory.
    \item \textbf{Multi-particle states.} These correspond to multiple excitations of the field, i.e. a collection of single-particle states, typically associated with a branch cut in the propagator of the theory.
    Such a state have a well-defined total momentum, but not a well-defined mass, since the internal energy of the state depends on the relative momenta of the single-particle states that compose it.
    \item \textbf{Bound states.} These are states that correspond to a bound state of multiple particles, again typically associated with a pole in the propagator of the theory.
    These states do have a well-defined momentum and mass, which essentially corresponds to the sum of the masses of the constituent particles minus the binding energy.
\end{itemize}
\end{mytheorem}